\documentclass{article}
\usepackage{amsmath, amssymb}
\title{Математическое обоснование GRA-Genius-Swarm}
\author{GRA Research Group}
\date{\today}

\begin{document}

\maketitle

\section{Формализация Интеллекта Роя}
Интеллект роя определяется как:
\begin{equation}
I_{swarm} = N \cdot \alpha \cdot C(\phi)
\end{equation}
где $C(\phi)$ — проектор согласованности, стремящийся к 1 при полном обнулении пены.

\section{Функционал Пены Мультиверса}
Для уровня $l$ пена определяется как:
\begin{equation}
\Phi^{(l)}(\Psi^{(l)}, G_l) = \sum_{\mathbf{a} \neq \mathbf{b}} \big| \langle \Psi^{(\mathbf{a})} | \mathcal{P}_{G_l} | \Psi^{(\mathbf{b})} \rangle \big|^2
\end{equation}
Цель алгоритма: найти состояние $\Psi^*$, такое что $\Phi^{(l)} = 0$ для всех $l$.

\section{Оценка Красоты (Интуиция)}
\begin{equation}
B(s) = \sigma\left( \sum w_i \cdot \text{sim}(h_i, \text{ideal}) \right)
\end{equation}
Решение считается гениальным, если $B(s) > 0.95$.

\section{Теорема о Полном Обнулении}
При выполнении условий коммутативности проекторов и полноты пространства, существует единственное состояние абсолютного когнитивного вакуума $\Psi_{\infty}^*$.

\end{document}